\chapter{Uso de inteligencia artificial como apoyo al desarrollo}
\label{appendix:ai-assisted-development}

Durante la realización del proyecto se utilizaron herramientas de inteligencia
artificial como apoyo puntual al trabajo del autor. Este uso debe distinguirse
de las capacidades de IA integradas en LIA, que constituyen una funcionalidad
de la aplicación y se describen en los capítulos de arquitectura e
implementación.

\section{Forma de uso}

La IA se utilizó como herramienta auxiliar para organizar el trabajo, dividir
objetivos en tareas y preparar listas de comprobación. También ayudó a detectar
errores, señalar incoherencias y sugerir aspectos que convenía revisar. El
autor valoraba cada observación, comprobaba el caso y decidía la actuación
correspondiente.

En la revisión técnica, la IA se empleó para orientar la búsqueda de causas
ante fallos de pruebas, tipos, compilación o integración, y para recordar las
comprobaciones que debían repetirse después de una modificación. En la memoria,
se empleó para revisar la organización de capítulos, detectar erratas y
contrastar la coherencia entre secciones, requisitos, implementación y
observaciones de la tutora.

Las instrucciones proporcionadas al asistente describían el objetivo de la
revisión, el contexto afectado, las restricciones y las comprobaciones
esperadas. Cuando una revisión era extensa, se separaba por bloques y el autor
contrastaba después sus resultados. Por ejemplo, ante los filtros de campos
personalizados se pidió revisar su coherencia entre listado, kanban y totales,
el aislamiento por cuenta y las pruebas relacionadas. En la memoria, las
instrucciones insistían en localizar contradicciones, no inventar evidencias,
distinguir lo realizado del trabajo futuro y comprobar la compilación.

\section{Responsabilidad y supervisión}

El autor mantuvo el control del alcance, los requisitos, las prioridades, la
arquitectura, las tecnologías y la aceptación final de cada cambio. Cada
observación se revisó y, cuando correspondía, se contrastó mediante pruebas
automatizadas, análisis de tipos, \emph{lint}, construcción de las aplicaciones
o compilación de la memoria. La responsabilidad sobre el contenido y el
resultado final del proyecto corresponde al autor.
